% =========================================================
\section{实验设计}
\label{sec:experiments}
% =========================================================

\subsection{时序运维 Benchmark}
\label{sec:dataset}

本文构建 TEF-RAG v6 Semantic Benchmark v1，用于受控评价动态运维场景中的时序证据检索。数据不是对真实电站故障日志的直接复制，而是在公开产品资料与公开研究数据给出的物理范围内，构造具有明确时间、因果与规程关系的运维事件链。其目标是系统控制晚到记录、版本替换、跨来源互补证据和持续不确定性等因素，并观察不同检索方法对这些因素的处理能力。

数据组织层次如表~\ref{tab:dataset_structure}。每条 chain 表示一个完整运维过程，由多条 evidence records 构成；每条 chain 派生多个不同任务意图，每个意图再生成语义保持的 paraphrase。检索 gold 以“必要证据组 + 必要 flow edges”表示，因此评价的不是某一条唯一文本，而是一个任务是否获得了足以形成完整流程的证据集合。

\begin{table}[t]
\caption{Benchmark 的数据层次}
\label{tab:dataset_structure}
\centering
\small
\begin{tabular}{@{}p{0.24\columnwidth}p{0.69\columnwidth}@{}}
\toprule
层次 & 含义 \\
\midrule
Chain & 一个完整运维事件过程 \\
Evidence & 可被检索的原子记录，带事件时间与可用时间 \\
Intent & 在同一 chain 上提出的不同运维信息需求 \\
Query & Intent 的语义保持改写 \\
Gold flow & 必要证据组及其关系约束 \\
\bottomrule
\end{tabular}
\end{table}

整个 benchmark 包含 400 条 authored chains、3888 条 evidence records、1200 个 primary intents 和 2400 条 query rows。开发、验证与测试按 chain 划分，从而避免同一事件过程的 paraphrase 跨 split。两类难度各包含 200 条 chain：Realistic Distribution Set 覆盖较常见的时序组合；Temporal-Hard Challenge Set 则强化晚到证据、supersession、跨 episode 混淆和跨来源互补证据。

\begin{table*}[t]
\caption{TEF-RAG v6 Semantic Benchmark v1 规模与划分}
\label{tab:dataset_stats}
\centering
\small
\begin{tabular}{lrrrr}
\toprule
统计项 & Development & Validation & Test & Total \\
\midrule
Authored chains & 240 & 80 & 80 & 400 \\
Primary intents & 720 & 240 & 240 & 1200 \\
Query rows & 1440 & 480 & 480 & 2400 \\
\midrule
Target assets & -- & -- & -- & 100 \\
Evidence records & -- & -- & -- & 3888 \\
\bottomrule
\end{tabular}
\end{table*}

公开资料只用于约束电芯类别、典型工作范围和时序采样背景，例如 HiTHIUM、REPT BATTERO 的储能电芯资料以及 RWTH M5BAT 数据~\cite{hithium2023v11,hithium2024v33,rept2025,rwth2026m5bat}。这些来源不用于定义本文合成的故障频率、站级阈值或具体维护因果故事。

% TODO: 增加一张数据结构图：Chain -> Evidence timeline -> Intents/Queries -> Required groups/flow edges。

\subsection{结构化生成评价}

为检验检索优势能否传递到下游任务，我们在同一批 authored chains 上建立结构化 generation gold。每个 primary intent 对应一个语义 gold object，并由其两个 paraphrases 共享。因此 development/validation/test 分别包含 720/240/240 个 semantic gold objects，对应 1440/480/480 条 queries。

每个输出由两部分组成：标准化工单 \texttt{work\_order} 与动作计划 \texttt{action\_plan}。工单描述诊断、推荐动作、适用规程、验证/不确定性状态及证据引用；动作计划由原子动作及其依赖关系构成。该评价关注“检索到的上下文是否足以支持更准确的结构化输出”，而不把语言模型本身作为本文的主要算法贡献。

\subsection{对比方法}
\label{sec:baselines}

正式比较五种检索上下文：BM25、BGE Reranker、Temporal-BM25、TA-RAG 与 TEF-RAG。BM25 代表词法检索；BGE Reranker 在 BM25 候选上进行通用语义重排；Temporal-BM25 在词法分数上加入时间感知项；TA-RAG 代表面向 diachronic QA 的时间感知 RAG；TEF-RAG 使用本文提出的关系图与集合级非线性排序。所有方法最终均输出不超过 Top-5 的 evidence IDs。

\subsection{评价指标}
\label{sec:metrics}

\textbf{Retrieval。}
使用 Recall@5 与 nDCG@5 衡量传统相关性；Complete@5 要求结果至少覆盖每个必要证据组；FlowComplete@5 在此基础上进一步要求存在全局一致的组--证据赋值，使必要 flow edges 同时得到满足。后两项直接衡量本文关注的集合级完整性。

\textbf{Generation。}
正文报告四项主要指标：Field Macro F1 (strict)、Action F1、Citation F1 与 Plan EM (strict)。它们分别衡量工单关键字段、规范化动作、claim--evidence 引用以及完整动作计划的严格匹配。其余 schema validity、dependency/order 与 end-to-end exact match 用于诊断。

\subsection{实验控制}

Retrieval 的开发、验证和最终测试严格分离；Generation comparison 固定使用同一个下游生成器、同一 prompt/schema 与解码设置，只改变输入的检索证据。这样可以把下游生成差异解释为检索上下文质量的变化，而不是生成模型配置差异。
